% !TEX root = 00_Template_SBSR.tex
\section{Conclusions}

The $\mathcal G^0_I$ distribution provided a better fit to the data than $\Gamma_{\text{SAR}}$ for both missions.
However, the Kullback--Leibler tests revealed differences among samples, mainly associated with mean brightness, indicating that a single point in the parameter space does not fully describe each class in the spatial domain.

For most classes and polarizations, the hypothesis of a common roughness parameter was not rejected, suggesting that distributions may share $\alpha$ while differing in $\mu$.
The exception was the NISAR Grassland class, for which this hypothesis was rejected in both HH and HV.
These findings support characterizing classes using families of $\mathcal G^0_I$ distributions rather than a single model with fixed parameters.

Future work should investigate which models best characterize the Grassland class for NISAR data; investigate how seasonality impacts the sample distributions over time, and assess whether distribution families improve land-cover classification across sensors and acquisition conditions.